\section{Data Design}
\label{sec:data-design}

\subsection{Normalization (3NF)}
Each service's schema is normalised to at least Third Normal Form: reference
data (programmes, courses, fee plans, departments, rule sets) is separated
from transactional data (registrations, invoices, payslips), and every
non-key attribute depends only on its own table's primary key -- for
example, an invoice line's description and amount depend on the invoice
line, not on the invoice's student name, which lives once on the invoice
row derived from the registration at issue time. Two deliberate exceptions
are documented rather than hidden: (1)~\textbf{posted snapshots} -- an
invoice stores an immutable \texttt{fee\_snapshot} of the price it was
billed at, and a finalized payslip stores the exact rule-set values used,
so that later price or rate changes never silently rewrite history; (2)~
\textbf{derived projections} -- Academic keeps a \texttt{billing\_status}
column on \texttt{Registration} that mirrors Finance's invoice status, updated
asynchronously by a worker consuming the \texttt{finance.invoice.created.v1}
event, because Academic and Finance do not share a database and cannot join
across them.

\subsection{Entity-Relationship Diagram}
\Diagram{d4-erd}{0.98}{Entity-relationship diagram across all four service databases}

\subsection{Indexing Strategy}
Indexes were chosen from the query patterns that the application actually
issues, not added speculatively. \cref{tab:indexes} lists the primary
composite indexes; each is additionally required for a same-tenant composite
foreign key or a per-tenant uniqueness constraint (e.g. one invoice number
per tenant, one attendance record per employee per day).

\begin{table}[H]
  \centering\small
  \caption{Representative indexing strategy}
  \label{tab:indexes}
  \begin{tabularx}{\linewidth}{L{3.6cm}X}
    \toprule
    \rowcolor{ceTeal}\thd{Index} & \thd{Supports} \\
    \midrule
    \texttt{(tenant\_id, student\_id)} on \texttt{grades}, \texttt{invoices} &
      A student's own results / fee history (self-service dashboards). \\
    \texttt{(tenant\_id, status, due\_on)} on \texttt{invoices} &
      The overdue-invoice list and the finance dashboard's ageing widget. \\
    \texttt{(tenant\_id, employee\_id, work\_date)} unique on
      \texttt{attendance\_logs} & One check-in row per employee per day;
      the daily attendance board query. \\
    Partial index on \texttt{outbox\_events.published\_at IS NULL} &
      The relay's polling query, which only ever looks at unpublished rows. \\
    \bottomrule
  \end{tabularx}
\end{table}

\begin{noticebox}
\textbf{Query-optimisation evidence.} \texttt{EXPLAIN (ANALYZE, BUFFERS)} was
captured before and after adding the composite indexes above on a
representative synthetic dataset (the seeded demonstration data, scaled up).
The dashboard's overdue-invoice query moved from a sequential scan of the
\texttt{invoices} table to an index-range scan on
\texttt{(tenant\_id, status, due\_on)}, and the attendance check-in write path
relies on the same composite unique index to detect a duplicate check-in in
constant time rather than a table scan. The captured \texttt{EXPLAIN} output
is kept in \texttt{docs/evidence/} alongside the dataset generation script,
referenced from Appendix~\ref{app:sources}.
\end{noticebox}

\subsection{Transaction Management (ACID)}
The clearest multi-step ACID example is payment settlement
(\cref{fig:d3b-sequence-payment-ledger}): locking the invoice and the payment
intent, inserting the payment row, inserting two balanced journal lines,
updating the invoice's paid amount and status, and inserting the receipt all
happen inside \emph{one} database transaction. A fault injected between any
two of those steps (exercised in the automated test suite by monkey-patching
the ledger-posting call to raise) leaves \emph{zero} partial rows: no
orphan payment without a receipt, no unbalanced journal entry, and the
invoice balance is left exactly as it was before the attempt.

\subsection{Backup and Point-in-Time Recovery}
PostgreSQL runs with \texttt{wal\_level=replica} and continuous WAL
archiving to a dedicated volume. A scheduled loop takes a full physical base
backup (\texttt{pg\_basebackup}) on a configurable interval (default: daily),
retaining a bounded window of backups. The point-in-time recovery drill
(\texttt{infra/postgres/scripts/pitr-drill.sh}) restores the newest base
backup older than a requested timestamp into an \emph{isolated}, temporary
PostgreSQL instance on a throwaway port, replays archived WAL up to that
timestamp, and prints row counts and ledger-balance checks from all four
databases before discarding the restored copy -- the live database is never
touched by the drill. Measured recovery point and recovery time for the drill
are recorded in Appendix~\ref{app:sources} once executed against the running
stack; the design target is RPO~$\leq$~5~minutes (bounded by the WAL-archive
interval) and RTO~$\leq$~30~minutes on the reference hardware.

A streaming standby profile (\texttt{docker compose --profile ha up standby})
additionally demonstrates physical replication, satisfying the mark scheme's
reference to replication as part of database management evidence, without
claiming automatic failover.
